% title: 试卷：复旦数学夏令营
% renderer: btex
% legacy-title: 用户:Solution/ 试卷: 复旦数学夏令营
% source: https://bananaspace.org/legacy/wiki/%E7%94%A8%E6%88%B7:Solution/_%E8%AF%95%E5%8D%B7:_%E5%A4%8D%E6%97%A6%E6%95%B0%E5%AD%A6%E5%A4%8F%E4%BB%A4%E8%90%A5
% summary: 从香蕉空间旧站迁移的原始 bTeX 页面。
\setdisplaytitle{试卷：学术营笔试}
\sethtmltitle{试卷：学术营笔试}

==2025年学术营笔试试卷==
\begin{enumerate}
    \item 设 $(V,\langle \cdot,\cdot\rangle)$ 为实内积空间，$T:V\to V$ 为线性映射。
    \begin{enumerate}
        \item 写出 $T$ 的伴随算子 $T^*$ 的定义。
        
        \item 考虑 $V=\mathbb{R}^2$ 及其一组基
        \[
        e_1=(1,0),\qquad e_2=(0,1),
        \]
        设 $V$ 上两个线性映射在这组基下的矩阵分别为
        \[
        A=
        \begin{pmatrix}
        3 & 1\\
        0 & -3
        \end{pmatrix},
        \qquad
        B=
        \begin{pmatrix}
        1 & 0\\
        -4 & -1
        \end{pmatrix}.
        \]
        若它们在 $V$ 的某个内积下均成为自伴随算子，求该内积在上述基下的矩阵。
    \end{enumerate}

    \item 求
    \[
    \int_0^{\frac{\pi}{2}} \left(\frac{x}{\sin x}\right)^2 \, dx .
    \]

    \item 求行列式
    \[
    D=
    \begin{vmatrix}
    1 & x_1 & \cdots & x_1^{\,n-2} & (x_2+\cdots+x_n)^{n-1}\\
    \vdots & \vdots & & \vdots & \vdots\\
    1 & x_n & \cdots & x_n^{\,n-2} & (x_1+\cdots+x_{n-1})^{n-1}
    \end{vmatrix}.
    \]

    \item 分析
    \[
    \int_0^{+\infty} \sin\left(x^p+\frac{1}{x^p}\right)\,dx
    \]
    在何时绝对收敛、条件收敛和发散。

    \item 设 $a_n$ 为方程
    \[
    x^n+nx-1=0
    \]
    的正实根，求级数
    \[
    \sum_{n=1}^{\infty} a_n \cos nx
    \]
    的收敛域 $D$，并问其在 $D$ 中是否绝对收敛和一致收敛。

    \item 设
    \[
    A=
    \begin{pmatrix}
    1 & 1 & 0\\
    0 & 1 & 0\\
    0 & 0 & 2
    \end{pmatrix}.
    \]
    \begin{enumerate}
        \item 问：$A$ 是否可对角化？
        
        \item 设
        \[
        \varphi:M(3,\mathbb{R})\to M(3,\mathbb{R}),\qquad
        X\mapsto A^{-1}XA,
        \]
        求 $\operatorname{tr}(\varphi^{100})$。
    \end{enumerate}

    \item 设
    \[
    \lim_{x\to+\infty}
    x\left(
    \int_0^x \frac{t}{\arctan t}\,dt
    -\frac{1}{\pi}x^2
    -\frac{4}{\pi^2}x
    -\frac{8}{\pi^3}\ln x
    +A
    \right)=B.
    \]
    求 $A$ 的表达式及 $B$ 的值。
\end{enumerate}